\section{Decay after clearing denominators}\label{sec:decay}
To prove that $D_nL_n\to0$, we compare denominator growth with decay
of the forms. The factorial growth of $D_n$ is matched by the factorial
decay of $L_n$. After these contributions cancel, we will bound the
remaining denominator growth rate by $8.2+15\log2$ and the decay rate
from below by $18.83$. Since
\[
 8.2+15\log2-18.83<-0.23,
\]
the cleared forms will tend to zero exponentially. The two denominator
terms come from the large odd primes and from $2$, respectively.
The small odd primes contribute only $o(n)$ after factorial normalization.

\subsection{The arithmetic envelope}
For a large prime write $x=n/p$. If $1<x<2$, a repeated address
class has its least representative between $0$ and $n-p$. Dividing
by $p$ restricts the corresponding variable $y$ to $[0,x-1]$.
For $x\ge2$, every residue class is repeated and we use a full period.

For real $x\ge1$, define
\begin{align*}
 N_j(x,y)&=\#\{h\in\Z:\delta_jx\le y+h\le(1-\delta_j)x\},\\
 C(x,y)&=\lfloor y+rx+\tfrac12\rfloor
                         -\lfloor y-(1+r)x+\tfrac12\rfloor,\\
 f(x,y)&=2+\sum_{j=0}^{30}N_j(x,y)-3C(x,y),\\
 J_x&=\begin{cases}[0,x-1],&1\le x<2,\\ {}[0,1],&x\ge2,\end{cases}
 \qquad
 e(x)=\max\left(0,\max_{y\in J_x}f(x,y)\right)-G\lfloor x\rfloor.
\end{align*}
The subtraction $G\lfloor x\rfloor$ removes the factorial contribution:
for $x=n/p$ and $p^2>7n$, Legendre's formula gives
$Gv_p(n!)=G\lfloor n/p\rfloor$.
At $x=n/p$ and $y=a/p$, $N_j(x,y)=N_{p,j}(a)$. Since $p$ is odd
and the interval endpoints are integers, $C(x,y)$ counts the
$\ell\in[-(1+r)n,rn-1]\cap\Z$ for which $p\mid2a+1+2\ell$.
Thus $3C(x,y)$ omits the two endpoint contributions and the prefactor indicator
from $c_p(a)$. Hence
Proposition~\ref{prop:isolated} gives
\begin{equation}\label{eq:prime-envelope}
 e_p-G\lfloor n/p\rfloor\le e(n/p).
\end{equation}
The restriction to $J_x$ for $x<2$ is where isolated pole classes
lower the estimate.

The interval defining $N_j$ is closed, so its exact expression is
\[
 N_j(x,y)=\lfloor y-\delta_jx\rfloor
               -\lceil y-(1-\delta_j)x\rceil+1.
\]
It is at most $(1-2\delta_j)x+1$, while $C(x,y)\ge(1+2r)x-1$.
Consequently
\begin{equation}\label{eq:envelope-tail}
 e(x)\le36+G\{x\},\qquad \{x\}=x-\lfloor x\rfloor.
\end{equation}
This bound also holds after replacing $e(x)$ by $\max(e(x),0)$.

\subsection{A finite partition and the prime sum}
To turn the prime costs into an exponential bound for $D_n$, we use
the prime number theorem. For fixed $1\le a<b$, it gives
\[
 \sum_{a<n/p<b}\log p=n(1/a-1/b)+o(n).
\]
The factor $1/a-1/b=\int_a^b dx/x^2$ explains why the total prime
cost is controlled by $\int_1^\infty e(x)\,dx/x^2$.
We compute the integral on $[1,400]$ by a finite partition and bound
the remaining tail.

The dependence on $y$ changes only at the affine offsets
\begin{equation}\label{eq:offsets}
 \delta_jx,\quad(1-\delta_j)x\quad(0\le j\le30),
 \qquad -rx-\tfrac12,\quad(1+r)x-\tfrac12,
\end{equation}
taken modulo one. On $[1,400]$, include every integer and every
point at which two distinct offsets coincide modulo one. For offsets
$sx+t$ and $s'x+t'$ with $s\ne s'$, these points are
$(t'-t+h)/(s-s')$, $h\in\Z$, whenever they lie in that interval.
There are $38\,372$ distinct event points.

Between consecutive events, the offsets keep their order, and
$\lfloor x\rfloor$ is unchanged. The moving boundary $x-1$ for
$x<2$ is itself the offset $x$ modulo one. Thus $e$ is constant on
each open event interval. At a non-event $x$, a breakpoint value
in $y$ equals one of its neighbouring cell values. At $y=0$ the
outer pole block takes its right-hand value; at $y=x-1$ it takes
its left-hand value. Checking the open cells inside $J_x$ therefore
also captures its boundary maxima.

For each event interval $(a,b)$, its contribution to
$\int_1^{400}e(x)\,dx/x^2$ is
$e((a+b)/2)(1/a-1/b)$. Exact rational computation, rounding each
contribution upwards to a multiple of $2^{-64}$, gives
\begin{equation}\label{eq:body-certificate}
 A_{\rm body}\le
 \frac{73794151590486879261}{9223372036854775808}.
\end{equation}
The rounding is upwards also for negative contributions.
For an integer $X\ge2$, \eqref{eq:envelope-tail} gives
\begin{align*}
 \int_X^\infty\frac{\max(e(x),0)}{x^2}\,dx
 &\le\frac{36}{X}+\frac G2\sum_{h=X}^\infty\frac1{h^2}\\
 &\le\frac{36}{X}+\frac{G}{2(X-1)}.
\end{align*}
At $X=400$ this is $34253/319200$. Together with
\eqref{eq:body-certificate}, the certified upper bound is
\begin{equation}\label{eq:A-certificate}
 A_{\rm cert}
 =\frac{745969417189431202995707}{92003136067626388684800}
 =8.10808684435616\ldots<\frac{41}{5}.
\end{equation}
The exact partition and two calculations of its cell maxima are
supplied with the paper; Section~\ref{sec:computations} describes them.

\begin{proposition}\label{prop:D-size}
As $n\to\infty$ through positive multiples of $96$,
\begin{equation}\label{eq:D-size}
 \log D_n\le G\log(n!)+n\left(\frac{41}{5}+15\log2+o(1)\right).
\end{equation}
\end{proposition}
\begin{proof}
The prime-number-theorem estimate above, summed over the finite
partition, yields its integral contribution.
Each event point can equal $n/p$ for at most one prime; by
\eqref{eq:envelope-tail}, all such primes together cost $O(\log n)$.

For the tail replace $e$ by $\max(e,0)$ on $[400,\infty)$.
The omitted small primes eventually lie there and may now be inserted
in an upper bound. On a fixed interval $[400,M]$, finite partitions
and the prime number theorem give the integral estimate. The remaining
primes $p\le n/M$ contribute at most
$(36+G)\vartheta(n/M)=O(n/M)$ by Chebyshev's bound, where
$\vartheta(u)=\sum_{p\le u}\log p$. Letting $M$ tend to infinity proves
\[
 \sum_{\substack{p\le n\\p^2>7n}}
     (e_p-G\lfloor n/p\rfloor)\log p
 \le n(A_{\rm cert}+o(1)).
\]
Add \eqref{eq:small-cost} and \eqref{eq:two-cost}. The small odd
primes cost $o(n)$, and $2$ costs at most $15n\log2+o(n)$.
Finally use $A_{\rm cert}<41/5$.
\end{proof}

\subsection{The real phase and its maximum}
Write $H(u)=u\log u$ for $u>0$, with $H(0)=0$, and set
\begin{equation}\label{eq:phase}
 \Phi(x)=3\bigl(H(x+1+r)-H(x-r)\bigr)
 -\sum_{j=0}^{30}\bigl(H(x+1-\delta_j)-H(x+\delta_j)\bigr),
 \quad x\ge r.
\end{equation}
This phase records the exponential size of the numerator and pole
products after their factorial normalization.

\begin{lemma}\label{lem:phase}
For $x\ge r$, $\Phi(x)<-18.83$. For $x\ge3$,
$\Phi'(x)\le-3/x$.
\end{lemma}
\begin{proof}
For $x>r$, the sign of $\Phi'(x)$ is the sign of
\[
 (x+1+r)^3\prod_{j=0}^{30}(x+\delta_j)
                  -(x-r)^3\prod_{j=0}^{30}(x+1-\delta_j).
\]
Put $x=r+z$ and clear positive denominators. The resulting polynomial
can be taken as
\begin{equation}\label{eq:critical-polynomial}
 P(z)=(2z+5)^3\prod_{j=0}^{30}(96z+72+j)
                       -8z^3\prod_{j=0}^{30}(96z+168-j).
\end{equation}
Its degree is $33$. Exact multiplication shows that its coefficients
in degrees $0,1,2$ are positive and all remaining coefficients are
negative. Descartes' rule of signs, together with $P(0)>0$ and its
negative leading coefficient, therefore gives exactly one positive
zero. Thus $\Phi$ increases to a unique global maximum and then
decreases.

Exact evaluations of $P$ bracket the maximum point $x_*$ by
\[
 \frac{378667493}{500000000}<x_*<\frac{757334987}{1000000000}.
\]
Rational logarithm enclosures on this interval give
\begin{equation}\label{eq:phase-certificate}
 \Phi(x_*)\le-\frac{18835112315041}{10^{12}}<-18.83.
\end{equation}
These finite checks are reproducible directly from
\eqref{eq:critical-polynomial}. For the logarithms, reduce the
argument to $1\le y\le2$ and use
\[
 \log y=2\sum_{j=0}^{47}\frac{w^{2j+1}}{2j+1}+R,
 \qquad w=\frac{y-1}{y+1},\qquad
 0\le R\le\frac{2w^{97}}{97(1-w^2)}.
\]
The same enclosure supplies $\log2$ when rescaling the argument.

To control the infinite tail, use $\log(1+u)\le u$ and
\[
 \log\frac{x+1-\delta_j}{x+\delta_j}
 =\int_{x+\delta_j}^{x+1-\delta_j}\frac{dt}{t}
 \ge\frac{1-2\delta_j}{x+1}.
\]
For $x\ge3$ these give
\[
 x\Phi'(x)\le3(1+2r)\frac{3}{3-r}-\frac{3\gamma}{4}
                  =-\frac{383}{64}<-3.
\]
\end{proof}

\subsection{Uniform estimates and the infinite sum}
For real $t\ge rn+1$, every factor of $R_n(t)$ is positive. Its
gamma-product expression is
\[
 R_n(t)=(2t+n+1)
 \left(\frac{\Gamma(t+(1+r)n+2)}{\Gamma(t-rn)}\right)^3
 \prod_{j=0}^{30}
 \frac{\Gamma(t+\delta_jn+\tfrac12)}
      {\Gamma(t+(1-\delta_j)n+\tfrac32)}.
\]
Stirling's formula, uniformly for real arguments at least one, yields
\begin{equation}\label{eq:stirling}
 \log R_n(t)=-G\log(n!)+n\Phi(t/n)+O(\log(n+t)).
\end{equation}
The potentially small argument $t-rn$ is still at least one, even
when $(t-rn)/n\to0$. Its Stirling remainder and the logarithmic
term omitted from $u\log u-u$ are therefore $O(\log(n+t))$.
Every other argument is bounded below by a positive multiple of $n$.
The $n\log n$ and linear terms give $-G\log(n!)$, up to $O(\log n)$.
This explains the uniformity of \eqref{eq:stirling} down to $t=rn+1$.

Differentiation costs only a power of $\log n$. Indeed,
\[
 \frac{R_n'(t)}{R_n(t)}
 =\frac{2}{2t+n+1}+3\sum_{\ell\in\Lambda_n}\frac1{t-\ell}
                       -\sum_{a=0}^{n}\frac{B_a}{t+a+\tfrac12}.
\]
The reciprocals of the consecutive zero distances sum to $O(\log n)$,
and their squares sum to $O(1)$, uniformly for $t\ge rn+1$.
The corresponding pole sums are bounded since $r>0$. Differentiating
this identity gives
\begin{equation}\label{eq:derivative-size}
 |R_n''(t)|\le C(1+\log n)^2R_n(t)
\end{equation}
for a fixed $C$.

Only terms $t\ge rn+1$ remain in $L_n$ by \eqref{eq:zeros}.
For $t\le3n$, equations \eqref{eq:stirling}--\eqref{eq:derivative-size}
and Lemma~\ref{lem:phase} bound the $O(n)$ terms by a polynomial
in $n$ times $(n!)^{-G}e^{-18.83n}$. For $t>3n$, the derivative
bound for $\Phi$ gives
\[
 \Phi(t/n)\le\Phi(3)-3\log(t/(3n)).
\]
The resulting power tail dominates the fixed factor $(1+t/n)^C$
from the error in \eqref{eq:stirling}; its sum adds only a polynomial
factor in $n$. We obtain
\begin{equation}\label{eq:L-size}
 |L_n|\le(n!)^{-G}\exp\bigl(n(-18.83+o(1))\bigr).
\end{equation}
Combining this with Proposition~\ref{prop:D-size},
\begin{equation}\label{eq:cleared-decay}
 \limsup_{\substack{n\to\infty\\96\mid n}}
       \frac1n\log|D_nL_n|
 \le\frac{41}{5}+15\log2-\frac{1883}{100}
 =-0.2327922916\ldots<-0.23.
\end{equation}
Thus the forms tend to zero after their coefficient denominators are
cleared. It remains to ensure that infinitely many are nonzero under
the rationality hypothesis.
